
\section{Round-twenty-seven reconstruction: quotient defect geometry and conditional cut cumulants}
\label{sec:r27-b3}

The collision-to-density map has an infinite-dimensional kernel.  The active
graph theorem therefore works on the Hilbert quotient of collision defects,
not on arbitrary representatives.  Process tightness is obtained from a
conditional cut-density theorem in the B2 history algebra; deterministic
restart from a complete microscopic point is never confused with stochastic
regeneration of an ensemble.

\subsection{The minimal collision-defect quotient}

Fix a positive regular Boltzmann path $f$ in the near-Maxwellian B2 chart and
let
\[
 A_f=\tfrac12ff_*B\,dt\,dx\,dv\,dv_*\,d\omega.
\]
Let
\[
 \mathcal H_c=L^2(A_f^{-1}),\qquad
 \mathcal C:\mathcal H_c\longrightarrow\mathcal Y
\]
be the closed gain-minus-loss map.  Put
\[
 \mathcal K=\ker\mathcal C,\qquad
 \mathcal H_c^\circ=\mathcal H_c/\mathcal K.
\]
Every class $[h]$ has the unique minimal representative
$h^\circ\in\mathcal K^\perp=\overline{\operatorname{Ran}\mathcal C^*}$ and
\[
 \|[h]\|_\circ=\|h^\circ\|_{\mathcal H_c}
 =\inf_{k\in\ker\mathcal C}\|h+k\|_{\mathcal H_c}.         \tag{B3.1}
\]
For a density tangent $u$ define the normal collision defect
\[
 h=\dot\Gamma-D A_f[u].
\]
The quotient balance operator is
\[
 \mathfrak B^\circ(u,[h])
 =\bigl(\partial_tu+v\cdot\nabla_xu-DQ_f[u]-\mathcal Ch^\circ,\ u(0)\bigr).
                                                                    \tag{B3.2}
\]

\begin{lemma}[No hidden collision kernel]
\label{lem:r27-b3-quotient}
If $\mathfrak B^\circ(0,[h])=0$, then $[h]=0$.  The collision part of the
quadratic action is
\[
 \tfrac12\|[h]\|_\circ^2,
\]
which is independent of the representative of the raw contact perturbation.
\end{lemma}

\begin{proof}
The first component gives $\mathcal Ch^\circ=0$.  Since
$h^\circ\in(\ker\mathcal C)^\perp$, it follows that $h^\circ=0$.
Orthogonal projection minimizes the Hilbert norm on each affine class, giving
the second statement.
\end{proof}

\subsection{Nonautonomous hypocoercive graph estimate}

Let $M$ be a Maxwellian and suppose
\[
 \|f/M-1\|_{W^{3,\infty}_{x,v}(M^{-1/2})}\le\epsilon_*    \tag{B3.3}
\]
uniformly on $[0,T]$.  Let $\Pi$ be the five-dimensional local collision
invariant projection.  The Maxwellian collision form gives
\[
 \|(I-\Pi)p\|_{L^2(\nu M)}^2
 \le C\int A_f(\Delta p)^2.                               \tag{B3.4}
\]
To control spatially varying coefficients, cover
$[0,T]\times\mathbb T^3$ by parabolic boxes of side $\ell$ and freeze the
coefficients at their centres.  Let $\chi_\alpha$ be a partition of unity and
let $K(D_x)$ be the standard order-zero Kawashima multiplier on nonzero
Fourier modes.

\begin{lemma}[Localized transport--collision observability]
\label{lem:r27-b3-observe}
For $\ell$ and $\epsilon_*$ sufficiently small, every adjoint test $p$ with
terminal condition $p(T)=0$ and global conservation gauge satisfies
\[
 \|p\|_{\mathcal X}^2
 \le C\left[
 \|-\partial_tp-v\cdot\nabla_xp-\mathcal L_f^*p\|_{\mathcal Y^*}^2
 +\int A_f(\Delta p)^2\right],                            \tag{B3.5}
\]
where $\mathcal X$ contains the microscopic
$L^2(\nu M)$ norm, the $H^{-1}_x$ norms of
$\partial_t\Pi p$ and $\nabla_x\Pi p$, and the endpoint traces.
\end{lemma}

\begin{proof}
For frozen Maxwellian coefficients, take spatial Fourier transform.  Estimate
(B3.4) controls the microscopic part.  The cross term
\[
 \operatorname{Re}\langle K(k)\widehat{\Pi p},
                         \widehat{(I-\Pi)p}\rangle
\]
controls the hydrodynamic mode for $k\ne0$; global conservation and the
terminal condition control $k=0$.  Localizing in $(t,x)$ creates commutators
$[\chi_\alpha,v\cdot\nabla_x]$ and
$[K(D_x),\chi_\alpha]$, bounded by $C\ell$ in the graph norm.  Coefficient
freezing contributes $C(\epsilon_*+\ell)$ by (B3.3).  Choose these constants
below one quarter of the frozen coercivity constant and absorb.  A finite
overlap sum gives (B3.5).  Thus spatial commutators, absent from a
time-only partition, are explicitly controlled.
\end{proof}

\begin{theorem}[Closed quotient balance range]
\label{thm:r27-b3-range}
After the five conservation conditions are imposed,
$\mathfrak B^\circ$ has closed range.  On the orthogonal complement of the
homogeneous linearized Boltzmann solutions it has a bounded right inverse
\[
 \mathcal R:\operatorname{Ran}\mathfrak B^\circ
 \longrightarrow\mathcal X_u\oplus\mathcal H_c^\circ.     \tag{B3.6}
\]
Its adjoint includes the initial-output component:
\[
 (\mathfrak B^\circ)^*(p,p_0)
 =\left(-\partial_tp-v\cdot\nabla_xp-\mathcal L_f^*p,\
        [\Delta p],\ p_0-p(0)\right),\quad p(T)=0.         \tag{B3.7}
\]
\end{theorem}

\begin{proof}
Lemma \ref{lem:r27-b3-observe} is the adjoint graph estimate modulo the
explicit finite-dimensional gauge.  The Hilbert closed-range theorem gives
closed range and a bounded inverse on the orthogonal complement.  Integration
by parts in time produces the terminal condition and the initial boundary
source $p_0-p(0)$; collision duality produces the class $[\Delta p]$ in the
quotient.  Lemma \ref{lem:r27-b3-quotient} removes the spurious
infinite-dimensional null directions.
\end{proof}

\subsection{A conditional cut-density theorem}

Let $\mathcal G_t^\varepsilon$ be the sigma-field generated by the empirical
density, contact current, and open-half-edge history state at time $t$, but not
by the unresolved microscopic labels inside future roots.  A bounded
$\mathcal G_t^\varepsilon$-stopping time is called a cut stopping time.

\begin{lemma}[Uniform conditional future density]
\label{lem:r27-b3-cut}
For every cut stopping time $\tau_\varepsilon$, the conditional law of the
future root variables on $[\tau_\varepsilon,\tau_\varepsilon+\delta]$ is
absolutely continuous with respect to the prepared B2 root reference, with
density $R_{\tau_\varepsilon}^\varepsilon$ satisfying
\[
 E\!\left[(R_{\tau_\varepsilon}^\varepsilon)^p
 \mid\mathcal G_{\tau_\varepsilon}^\varepsilon\right]\le C_p            \tag{B3.8}
\]
for some $p>1$, uniformly on the regular history sublevel.  Conditional
connected cumulants obey
\[
 |\kappa_r^\tau(\Theta_1,\ldots,\Theta_r)|
 \le C_r\mu_\varepsilon^{1-r/2}
       (\delta+\mu_\varepsilon^{-1})\prod_j\|\Theta_j\|_{\mathcal S_m}.
                                                                    \tag{B3.9}
\]
\end{lemma}

\begin{proof}
Cut every B2 history at $\tau_\varepsilon$.  The open half-edges are
$\mathcal G_\tau^\varepsilon$-measurable and the unresolved future roots retain
the grand-canonical factorial reference.  The radius-loss estimate
(B2.11), applied simultaneously to the numerator and denominator of the
conditional density, gives an $L^p$ bound after shrinking the history radius.
Connected future graphs must contain a first vertex after the cut; its time
integral gives $\delta$, and a single microscopic cell gives the additional
$\mu_\varepsilon^{-1}$.  The automorphism normalization makes the remaining
genealogy sum convergent.  This proves (B3.8)--(B3.9).  Conditioning on the
complete phase point would instead give a Dirac law and is not used.
\end{proof}

\subsection{Nuclear process tightness}

Choose Fourier modes in $(t,x)$, Hermite modes in $v$, and spherical harmonics
in $\omega$.  Let $\mathcal S_m$ have coefficient weight
$(1+|k|+|n|+|\ell|)^m$ and Gaussian velocity weight.  For a rapidly increasing
subsequence, $\mathcal S_{m_{j+1}}\hookrightarrow\mathcal S_{m_j}$ is
Hilbert--Schmidt; hence
\[
 \mathcal S=\bigcap_j\mathcal S_{m_j}
 \subset\mathcal H\subset
 \mathcal S'=\bigcup_j\mathcal S_{-m_j}
\]
is nuclear.

Let $Z^\varepsilon$ be the centred density/contact fluctuation process.

\begin{lemma}[Aldous estimate from cut cumulants]
\label{lem:r27-b3-aldous}
For every cut stopping time and $0\le\theta\le\delta$,
\[
 E\|Z^\varepsilon_{\tau_\varepsilon+\theta}
      -Z^\varepsilon_{\tau_\varepsilon}\|_{\mathcal S_{-m}}^2
 \le C(\delta+\mu_\varepsilon^{-1}),                      \tag{B3.10}
\]
and the maximum jump tends to zero in probability.
\end{lemma}

\begin{proof}
Apply (B3.9) to an orthonormal basis of $\mathcal S_m$ and sum using the
Hilbert--Schmidt embedding.  The fourth-cumulant bound gives compact
containment.  A single empirical atom has
$\mathcal S_{-m}$ norm $O(\mu_\varepsilon^{-1/2})$; B2 factorial contact
tails control simultaneous jumps.  Aldous' criterion follows.
\end{proof}

\begin{theorem}[Joint kinetic process CLT]
\label{thm:r27-b3-clt}
The density/contact fluctuations converge in
$D([0,T];\mathcal S')$ to the continuous Gaussian solution of the quotient
linearized balance equation.  Its collision-noise covariance is
\[
 E[dW_c(\phi)dW_c(\psi)]
 =\int A_f\,\Delta\phi\,\Delta\psi,                        \tag{B3.11}
\]
and exact-number preparation replaces the initial covariance by the B1 Schur
complement.
\end{theorem}

\begin{proof}
B2 pressure derivatives identify the second cumulant, while higher connected
cumulants vanish after central-limit scaling.  Lemma
\ref{lem:r27-b3-aldous} and nuclear compact containment give process
tightness.  The exact microscopic Green identity identifies the limiting
martingale problem.  Uniqueness follows from the adjoint observability
estimate (B3.5).
\end{proof}

\subsection{The exact second epi-derivative}

On a zero-cost path $\Gamma=A_f$, use the normal quotient coordinate
$[h]=[\dot\Gamma-D A_f[u]]$.  A positive exponential chart is
\[
 f_\eta=f\exp\{\eta u/f-\psi_\eta\},\qquad
 \Gamma_\eta=A_{f_\eta}
 \exp\{\eta h^\circ/A_f+r_\eta\}.                         \tag{B3.12}
\]
The first-order balance residual is zero.  The order-$\eta^2$ residual lies
in the compatibility range of Theorem \ref{thm:r27-b3-range}; applying
$\mathcal R$ determines $r_\eta$ and a density correction.  Truncation is
performed before the correction and removed diagonally.

\begin{theorem}[Quotient Mosco Hessian]
\label{thm:r27-b3-mosco}
At $(f,A_f)$,
\[
 d_e^2I[u,[h]]
 =\tfrac12\|u_0\|_{\Sigma_0^{-1}}^2
  +\tfrac12\|[h]\|_\circ^2,                               \tag{B3.13}
\]
subject to the quotient linearized balance and microcanonical tangent
conditions, and $+\infty$ otherwise.  The forms Mosco-converge and their
inverse on the closed range equals the covariance of Theorem
\ref{thm:r27-b3-clt}.
\end{theorem}

\begin{proof}
For the lower bound, bounded rescaled entropy gives a de la
Vall\'ee--Poussin bound for the normal defects, hence uniform integrability.
Weak limits satisfy the quotient balance equation and convex lower
semicontinuity gives (B3.13).  For the upper bound, use (B3.12) and the bounded
right inverse; the correction is $O(\eta^2)$ in graph norm and does not alter
the quadratic cost.  Graph coercivity gives equicoercivity.  The
Cameron--Martin minimization of (B3.11) over collision noises is exactly the
minimal quotient norm, proving the covariance identification.
\end{proof}

This quotient construction eliminates the explicit $\ker\mathcal C$
counterexample, while the cut-density theorem replaces the invalid full-state
restart and supplies the process estimate needed downstream.
